\section{Выводы}
\begin{enumerate}
  \item Алгоритм Ахо~--~Корасик можно адаптировать для поиска с
        джокерами, разбивая образец на твёрдые фрагменты, находя их
        автоматом за линейное время и проверяя стыковку через подсчёт
        голосов.

  \item Сложность алгоритма составляет $O(n + m + Z)$, где $Z$~---
        количество вхождений фрагментов в текст. Когда фрагменты
        длинные и редкие, $Z$ мало и алгоритм работает за линейное
        время, обходя наивный перебор (в тестах~--- в $\sim$1{,}8
        раза быстрее). Если же фрагменты короткие и часто
        встречаются, $Z$ достигает $O(n \cdot m)$, и продвинутый
        алгоритм может проигрывать наивному (в тестах~--- в $\sim$2{,}2
        раза медленнее). Таким образом, выбор алгоритма должен
        опираться на структуру образца, а не на формальную
        «продвинутость» метода.
\end{enumerate}
\pagebreak